\section{Conclusion}
\label{sec:conclusion}

We asked whether a jointly learned graph improves multivariate stock-return forecasting compared with conventional time-series models and a derivative-free graph search. On 20 semiconductor and agriculture companies over 1{,}806 trading days, under a chronological leakage-controlled protocol with three seeds, the answer is no on all counts.

\paragraph{Answers to the research question.}
\begin{itemize}[nosep]
    \item \textbf{Against non-graph models:} MTGNN with a learned graph attains test MAE $0.020948 \pm 0.000018$, worse than a zero forecast ($0.020869$) and worse than VAR ($0.020820$). Its information coefficient is $-0.0081$ with a 95\% interval of $[-0.050, +0.034]$. A Diebold-Mariano test against the zero forecast gives $p = 0.70$.
    \item \textbf{Against derivative-free search:} CMA-ES, searching the identical graph parameterization without gradients, attains $0.020951 \pm 0.000037$, statistically indistinguishable from the gradient-learned graph ($p = 0.85$), while costing roughly 11 minutes of search per seed on top of a pretrained backbone and a final fit. The two tie on accuracy; gradient learning wins decisively on cost.
    \item \textbf{Does the graph help at all:} No. Learned, correlation-based, CMA-ES-searched and identity adjacencies all fall within $2.0 \times 10^{-5}$ MAE of each other, with the identity, which disables message passing entirely, nominally best.
\end{itemize}

Across fourteen Diebold-Mariano comparisons, only two reach significance at the 5\% level, and in both the model is significantly \emph{worse} than predicting zero. No model's IC confidence interval excludes zero.

\paragraph{Why.} Two mechanisms account for every result. First, \textbf{shrinkage}: on a target that is close to unpredictable, the MSE-optimal forecast is the conditional mean, which is approximately zero, so every trained model collapsed its prediction dispersion to $5$--$18\%$ of the target's. This makes MAE reward hedging rather than skill, which is why the near-constant forecasts top the error table and why persistence, the only method predicting at full scale, sits last. Second, \textbf{a flat objective}: graph convolution routes information between nodes, and with no cross-sectional signal to route, the loss is insensitive to the adjacency matrix. The gradient reaching the node embeddings is therefore dominated by noise, which is why the learned graph is close to uninterpretable ($51.2\%$ within-sector against a $47.4\%$ random baseline) while a fixed correlation graph recovers coherent economic clusters ($89.8\%$) and still forecasts no better.

\paragraph{Contributions.} We transferred graph structure learning from its usual low-noise benchmarks to financial returns and reported the outcome under a protocol we verified rather than assumed, including automated tests for split leakage. We provided the gradient-versus-evolutionary graph-search comparison that we believe had not previously been made for this task, measuring cost as well as accuracy. Most usefully, we showed that the standard metrics for this task are dominated by prediction dispersion rather than forecast quality, and we documented three concrete ways this study could have reported a false positive: ranking by MAE, where shrinkage looks like skill; VAR's directional accuracy of $0.5332$, which came from an AIC-selected zero-lag model whose predictions are bit-identical to a constant; and the zero forecast's net Sharpe ratio of $1.694$, which came from an \texttt{argsort} tie-break over a constant basket. Each was caught only because an uninformative reference sat beside it in the same table.

\paragraph{Limitations.} The usable window began in May 2019 because Nutrien and Corteva did not exist earlier, costing roughly four years of history for the other 18 companies; we chose to preserve the specified company universe over sample length. The test period is a single contiguous stretch of 271 days rather than a rolling-origin evaluation. During the CMA-ES search the backbone was frozen rather than co-adapted, so our cost-accuracy comparison holds under a realistic budget rather than an unlimited one. We did not tune hyperparameters extensively, deliberately: in this signal regime, extended search on validation data invites selecting noise \parencite{harvey2016cross}.

\paragraph{Future work.} The most promising direction is changing the target rather than the model. Volatility and realised correlation are markedly more persistent than returns and are intrinsically relational, so a graph may earn its place predicting tomorrow's covariance structure even though it cannot predict tomorrow's return. Longer horizons are similarly more plausible, since fundamental linkages need time to propagate along a supply chain, whereas a one-day horizon is dominated by microstructure noise. A time-varying or regime-conditional graph would address the non-stationarity suggested by the decay of directional accuracy from $0.5307$ to $0.4959$ across volatility quartiles. Finally, a rolling-origin evaluation across several market regimes would test whether these conclusions are window-specific, though we note the ablation conclusion cannot be, since all graph variants were evaluated on identical data and were indistinguishable.

\paragraph{Closing remark.} The honest summary is that graph structure learning did not help here, and that we can explain precisely why: the method routes information that this target does not contain. That is a result about the limits of the approach rather than a failure to apply it, and it is only visible because we measured against references that contain no information by construction.
